%% kapa: exam=1 round=34 year=2023 period=2 form=A subject=accounting subject_ko=회계학 seq=8 number=48
%% source: Q-Net 감정평가사 기출문제 (cst00309) artlSeq=5211789 file=raw/1008/5211789/2023년도 제34회 감정평가사 1차시험 2교시.pdf
%% conversion: kordoc 4.24.51 -> tools/gen1.py (auto); review=proofread-sonnet; flags=none
\begin{question}{48}
퇴직급여제도의 용어에 관한 설명으로 옳은 것은?
\begin{choices}
  \item 비가득급여: 종업원의 미래 계속 근무와 관계없이 퇴직급여제도에 따라 받을 권리가 있는 급여
  \item 약정퇴직급여의 보험수리적 현재가치: 퇴직급여제도에 의거하여 현직 및 전직 종업원에게 이미 제공한 근무용역에 대해 지급할 예상퇴직급여의 현재가치
  \item 급여지급에 이용가능한 총자산: 제도의 자산에서 약정퇴직급여의 보험수리적 현재가치를 제외한 부채를 차감한 잔액
  \item 확정기여제도: 종업원에게 지급할 퇴직급여금액이 일반적으로 종업원의 임금과 근무연수에 기초하는 산정식에 의해 결정되는 퇴직급여제도
  \item 기금적립: 퇴직급여를 지급할 현재의무를 충족하기 위해 사용자와는 구별된 실체(기금)에 자산을 이전하는 것
\end{choices}
\end{question}
\answer{2}
